\chapter{统计推断与随机实验}
\label{ch:week06-ab-testing}

两组转化率的差异可能来自处理效果、对象构成或随机波动。随机实验通过事先规定
处理分配机制，为组间比较提供明确的概率依据。效应估计及其精度取决于实验
单元、潜在结果和目标参数；是否推广还须结合实施成本与风险判断。

随机化为组间效应比较提供概率依据，方差描述估计精度，功效用于确定样本量，
风险指标和经济阈值限定推广条件。协变量调整利用实验前信息，整群设计以共同
接受干预的对象为分配单位，置信序列控制持续观察中的覆盖错误。各方法的有效性
取决于相应的分配与观测条件。

\section{经营比较与统计推断}

描述性比较给出样本中的差异。统计推断进一步研究这种差异对目标总体的含义及其抽样不确定性；随机实验则通过已知分配机制研究处理效应。实施推广时，需要将效应估计与成本和风险条件相联系。

描述统计面对的是已经观察到的数据。统计推断研究未完全观察的总体，并利用重复抽样机制刻画样本结论的不确定性。若 \(Y\) 表示某个管理指标，例如是否转化、订单收入或页面加载时间，样本均值
\[
  \bar Y=\frac{1}{n}\sum_{i=1}^n Y_i
\]
只是总体均值 \(\mu=\E[Y]\) 的估计。不同样本会产生不同的 \(\bar Y\)。因此，管理报告不能只写“B 组转化率比 A 组高 20 个百分点”，还要说明这种估计的不确定性、样本量、适用对象和业务阈值。

\begin{definition}[统计推断]
统计推断是利用样本数据对未完全观察的总体参数或未来重复样本行为作出带不确定性的判断。它关注估计量、标准误、置信区间、检验和决策风险。
\end{definition}

统计推断与因果评估的目标有所不同：前者可以描述样本差异的不确定性，后者还需要判断差异是否由某项行动造成。随机化减少选择偏差，并使差异估计量的抽样不确定性具有可分析的概率依据。

\section{随机实验与统计推断}
\subsection{总体、样本、估计量与标准误}
\subjectindex{biaozhunwu}{标准误}

总体是研究所针对的对象集合或产生观测的概率分布。例如，本季度全部订单构成一个有限总体，未来访客的访问行为则可用概率模型描述。样本是按照给定机制取得的观测。参数是总体的数量特征，估计量是用于估计参数的样本统计量。有限总体推断与概率模型下的推断都需明确重复抽样或随机分配的机制。

\begin{definition}[标准误]
设 \(\hat\theta\) 是总体参数 \(\theta\) 的估计量。标准误是估计量在重复抽样下的标准差，记为
\[
  \operatorname{SE}(\hat\theta)=\sqrt{\operatorname{Var}(\hat\theta)}.
\]
在实际数据中，标准误通常需要用样本方差或模型估计来近似。
\end{definition}

对独立同分布且具有有限方差的观测，若样本方差为 \(s^2\)，则均值标准误的常用估计为
\[
  \widehat{\operatorname{SE}}(\bar Y)=\frac{s}{\sqrt n}.
\]
在这些条件下，均值标准误按 \(1/\sqrt n\) 的速度随样本量下降。若希望把标准误减半，通常需要约四倍样本量。样本量不足时，较大的观测差异也可能具有较宽的置信区间，不足以支持推广。

\subsection{随机实验的基本要素}

A/B 测试是一种随机对照实验。它把符合条件的实验单元随机分配到对照组和处理组，再比较结果。在适当随机分配和无干扰条件下，组间平均差异可无偏估计实验对象的平均处理效应。两组在某次实验中仍可存在随机的基线差异，这部分不确定性由分配机制刻画。

\begin{definition}[A/B 测试的基本元素]
一个面向管理决策的 A/B 测试至少包含：
\begin{enumerate}
  \item 实验单元：被随机分配的对象，如访客、订单、门店、销售人员或客户；
  \item 处理与对照：\(T_i=1\) 表示处理组，\(T_i=0\) 表示对照组；
  \item 主指标：用于判断实验是否达到核心目标的指标；
  \item 约束指标：用于防止局部优化损害系统安全、体验或长期价值的指标；
  \item 分析窗口：从分配到观察结果的时间范围；
  \item 纳入和排除规则：哪些对象有资格进入实验，哪些记录应被排除；
  \item 决策规则：何时推广、继续实验、停止实验或回滚。
\end{enumerate}
\end{definition}

A/B 测试可用于页面、推荐、优惠、履约承诺或运营流程的比较。例如，比较两种结账页面时，以页面版本为处理，以分配后规定时限内是否完成购买为结果。

\subsection{潜在结果与随机化}
\subjectindex{qianzaijieguo}{潜在结果}
\subjectindex{SUTVA}{SUTVA}
\subjectindex{suijishiyan}{随机实验}

潜在结果记号用于描述同一对象在不同处理下的结果。假定处理版本明确，且一个单元的结果不受其他单元处理分配的影响。这两项通常合称稳定单元处理值假设（stable unit treatment value assumption，SUTVA）。对实验单元 \(i\)，令 \(Y_i(1)\) 表示接受处理时的结果，\(Y_i(0)\) 表示接受对照时的结果。个体处理效应为
\[
  \tau_i=Y_i(1)-Y_i(0).
\]
但同一个实验单元不能同时处于处理和对照状态，所以 \(\tau_i\) 无法被直接观察。实验目标通常是平均处理效应：
\[
  \tau=\E[Y(1)-Y(0)].
\]

若处理分配 \(T_i\) 是随机的，且实验执行没有系统性污染，那么处理组均值与对照组均值之差
\[
  \hat\tau
  =\bar Y_1-\bar Y_0
  =\frac{1}{n_1}\sum_{i:T_i=1}Y_i
   -\frac{1}{n_0}\sum_{i:T_i=0}Y_i
\]
就是 \(\tau\) 的自然估计量。

\begin{proposition}[随机化差异估计量的无偏性]
在固定 \(N\) 个实验单元中，采用完全随机分配，使 \(n_1>0\) 个单元接受处理、\(n_0=N-n_1>0\) 个单元接受对照；假定处理版本明确、无单元间干扰、结果完整观测。对随机分配取期望，则
\[
  \E_T[\bar Y_1-\bar Y_0]=\frac1N\sum_{i=1}^N[Y_i(1)-Y_i(0)].
\]
\end{proposition}

证明利用每个单元具有相同的入组概率。完全随机分配下，\(\E_T[T_i]=n_1/N\)，
于是
\[
\E_T\!\left[\frac{1}{n_1}\sum_iT_iY_i(1)\right]
=\frac{1}{N}\sum_iY_i(1).
\]
对照组同理，两式相减即得结论。这里的期望针对处理分配，实验对象的潜在结果视为固定。
推广到其他人群时，还须说明实验样本与目标总体的关系。

\subsection{随机化估计量的方差}
\subjectindex{Neyman}{Neyman方差}
在上述有限总体随机化框架中，令 \(S_1^2,S_0^2\) 为两组潜在结果在全部 \(N\) 个单元中的有限总体方差，\(S_\tau^2\) 为个体处理效应的方差，分母均使用 \(N-1\)。Neyman 方差为
\[
\operatorname{Var}_T(\hat\tau)
=\frac{S_1^2}{n_1}+\frac{S_0^2}{n_0}-\frac{S_\tau^2}{N}.
\]
最后一项无法直接估计，因为同一个体的两个潜在结果不能同时观察。去掉非负的 \(S_\tau^2/N\)，在 \(n_1,n_0\ge2\) 时，用组内样本方差估计前两项，得到常用的保守方差估计。在独立同分布的总体抽样框架下，也可得到两组均值差的标准误公式；此时的期望与方差针对总体抽样，而此处的随机性仅来自处理分配，两者应加以区分。

随机化无偏性是同一设计反复实施时的平均性质，并不保证每次实验中的两组都完全平衡。若处理组出现更多缺失、使用处理后的资格条件删样本，实际分析就偏离了上述命题的前提。

\section{区间、功效与决策阈值}
\label{sec:inference-interpretation}
\subsection{均值差异、置信区间与 p 值}
\subjectindex{zhixinqujian}{置信区间}
\subjectindex{pvalue}{p值}

对连续指标或二元指标，都可以从均值差异开始。二元转化指标 \(Y_i\in\{0,1\}\) 的均值就是转化率。设处理组和对照组样本均值分别为 \(\bar Y_1,\bar Y_0\)，样本方差分别为 \(s_1^2,s_0^2\)，则差异估计量的标准误可写为
\[
  \widehat{\operatorname{SE}}(\hat\tau)
  =
  \sqrt{\frac{s_1^2}{n_1}+\frac{s_0^2}{n_0}}.
\]
在适用的中心极限定理条件下，令 \(z_p\) 为标准正态分布的 \(p\) 分位数，
\(100(1-\alpha)\%\) 近似置信区间为
\[
  \hat\tau
  \pm
  z_{1-\alpha/2}\widehat{\operatorname{SE}}(\hat\tau).
\]

\begin{definition}[p 值]
在给定原假设 \(H_0\) 的条件下，p 值是观察到当前统计量或更极端统计量的概率。概率计算针对原假设规定的统计量分布。
\end{definition}

检验回答样本结果与给定原假设是否相容，推广决策还需比较效果大小与实施成本。
因此，报告应同时列出效应估计、区间、检验结果和预先规定的业务阈值。区间较宽时，
是否追加实验取决于预期信息价值和新的实验设计。

\subsection{置信区间的解释}
频率学派将参数视为固定值，区间的随机性来自样本。若区间构造规则记为
\(C(\mathcal D)\)，95\%置信水平表示在规定的重复抽样机制下，
\(\Prob_\theta\{\theta\in C(\mathcal D)\}\) 达到相应覆盖水平。
贝叶斯可信区间则根据先验和后验分布描述参数的不确定性。两种区间的概率解释各有依据。

\begin{example}[幅度、精度与经济阈值]
设两组相互独立，每组400个二元观测，对照组转化率为0.35，处理组为0.40。
采用大样本未合并标准误，得到
\[
\widehat{\rm SE}=
\sqrt{\frac{0.35(1-0.35)}{400}+\frac{0.40(1-0.40)}{400}}
\approx0.0342.
\]
均值差为0.05，近似95\%区间为
\(0.05\pm1.96\times0.0342\)，即约 \([-0.017,0.117]\)。
效应的点估计为正，但按当前估计精度，小幅负效应与较大正效应都不能排除。
如果至少提升0.03才值得推广，区间也没有排除低于业务阈值的效应。
\end{example}

若研究问题要求验证新版本的损失处于可接受范围，应预先规定非劣效界限；若要求两版本足够接近，应规定等效区间并采用等效性检验。差异检验、非劣效性检验和等效性检验具有不同的原假设，实验设计应由研究问题确定。

继续收集数据会增加成本，并可能伴随用户与环境变化。实验设计应预先规定目标效应、精度、停止规则及最大样本量。若研究过程中修改方案，应记录修改时点、原因及相应的推断方法。

\subsection{显著性、功效与最小可检测效应}
\subjectindex{gongxiao}{功效}
\subjectindex{MDE}{MDE（最小可检测效应）}

样本量由目标效应、结果变异和容许错误率共同决定。显著性水平 \(\alpha\) 约束原假设成立时错误拒绝它的概率；功效 \(1-\beta\) 是在指定备择效应下拒绝原假设的概率。检验控制统计错误率，推广决策还需满足净收益和风险约束。

最小可检测效应，通常记为 MDE，是在给定检验、显著性水平、样本量与目标功效下能够达到该功效的效应大小。它是统计设计的灵敏度，不自动等于最小值得实施的业务效应。例如，当前转化率为 35\% 时，组织可能希望对 10 个百分点的提升具有 80\% 的检出功效，同时另行依据成本确定推广门槛。

对两组比例差异，一个常用的每组样本量近似为
\[
n
\approx
\frac{
\left[
z_{1-\alpha/2}\sqrt{2\bar p(1-\bar p)}
+
z_{1-\beta}\sqrt{p_0(1-p_0)+p_1(1-p_1)}
\right]^2
}{
(p_1-p_0)^2
},
\]
其中 \(p_0\) 是基准转化率，\(p_1=p_0+\delta\)，\(\delta\) 是 MDE，\(\bar p=(p_0+p_1)/2\)，\(1-\beta\) 是检验功效。

\begin{modernbox}[管理解释]
应分别报告统计设计的 MDE 与业务上的最低净收益门槛。若希望检出的差异很小，所需样本量可能超出可承受范围；若 MDE 远高于具有业务价值的差异，实验可能对重要改进缺乏功效。设计阶段需要统筹精度、成本和实施期限，使统计目标与推广条件相匹配。
\end{modernbox}

\subsection{约束指标与实验决策}
\subjectindex{yueshuzhibiao}{约束指标}

主指标改善可能伴随其他结果恶化。例如，新版结账页提高转化率的同时增加加载时间，优惠规则提高订单量的同时减少毛利。实验因而需要预先确定风险指标和容许变化范围。

\begin{definition}[约束指标]
约束指标（guardrail metric）用于限制实验推广可能带来的不利影响。研究者事先规定各指标的可接受范围，推广方案须同时满足主指标的收益要求和这些风险限制。
\end{definition}

约束指标常见于以下几类：
\begin{table}[htbp]
\centering
\caption{A/B 测试中的常见约束指标}
\begin{tabular}{P{0.20\linewidth}P{0.34\linewidth}P{0.34\linewidth}}
\toprule
\rowcolor{BDMTableHead}
类型 & 指标例子 & 管理含义\\
\midrule

体验约束 & 页面加载时间、投诉率、客服工单率 & 避免短期转化提升损害用户体验。\\
财务约束 & 毛利率、补贴成本、退货率 & 避免收入提升来自不可持续补贴或高退货。\\
运营约束 & 履约延迟、库存缺货、人工处理量 & 避免前端策略把压力转移到供应链或运营团队。\\
长期约束 & 复购率、留存率、订阅取消率 & 避免短期行为伤害长期客户价值。\\
\bottomrule
\end{tabular}
\end{table}

实验报告应同时呈现主指标与风险指标的效应估计、置信区间、业务阈值和样本量。
推广须满足预设收益条件及风险容限。未满足条件时，可结束当前试验；若拟修改方案
或补充数据，应另行明确设计和停止规则。探索性分层用于提出后续研究假设。

\section{协变量调整与制造实验设计}

\subsection{用实验前协变量降低方差}
\subjectindex{CUPED}{CUPED}
实验前已经观察到的订单量、客户活跃度或设备历史负载，可以解释结果的部分变化。令 \(X\) 是实验前变量，构造调整后的结果
\[
Y^{\rm adj}=Y-\theta(X-\mu_X).
\]
若随机化使两组的 \(\E[X]\) 相同，且两组使用共同的 \(\theta,\mu_X\)，调整不会改变总体处理效应。它可以减少与处理无关的基线波动。展开方差：
\[
\operatorname{Var}(Y^{\rm adj})
=\operatorname{Var}(Y)+\theta^2\operatorname{Var}(X)
-2\theta\operatorname{Cov}(X,Y).
\]
当 \(\operatorname{Var}(X)>0\) 时，最小化得到
\[
\theta^*=\frac{\operatorname{Cov}(X,Y)}{\operatorname{Var}(X)},
\qquad
\operatorname{Var}(Y^{\rm adj})=\operatorname{Var}(Y)(1-\rho_{XY}^2).
\]
这就是 CUPED\footnote{CUPED 是 Controlled-experiment Using Pre-Experiment Data 的缩写，指利用实验前数据降低随机实验估计方差的方法。} 所利用的基本方差缩减思想\cite{cuped06}。上述方差公式在给定总体最优系数时精确成立；有限样本估计系数需要考虑估计误差，可用独立历史数据或交叉拟合\footnote{交叉拟合将样本分组，用其他组的数据估计调整模型，再对留出组计算调整结果，以避免同一观测同时用于拟合和评价。第8章进一步讨论其在因果估计中的作用。}。不得使用会被处理影响的实验后变量来调整。

\begin{example}[方差缩减的数量级]
若基线与结果相关系数为 \(0.6\)，且采用总体最优调整系数，则方差比例为 \(1-0.36=0.64\)，标准误比例为 \(0.8\)。在效应和设计保持相同的大样本近似下，所需样本量约为原来的 \(64\%\)。
\end{example}

\subsection{制造实验的随机化单元}
\subjectindex{zhengqunsheji}{整群设计}
\subjectindex{jiaochaganrao}{干扰}
\subjectindex{suijihua_danyuan}{随机化单元}
维护策略的随机分配单位通常是整台设备。同一设备产出的产品共享温度、班组和故障冲击，不能将数万件产品视为数万个独立实验单元。

若每个集群有 \(m\) 个观测、集群内相关系数为 \(\rho_c\)，在等规模且交换相关的简化模型下，均值方差相对独立样本放大
\[
\operatorname{DE}=1+(m-1)\rho_c.
\]
当 \(m=50,\rho_c=0.1\) 时，设计效应为 \(5.9\)。增加集群内件数并不能替代增加独立设备或班组。若处理按设备随机，应在设备层面分析或使用适当的集群稳健推断；只有极少数集群时，普通渐近近似仍不可靠。

\begin{figure}[htbp]
\centering
\begin{tikzpicture}[node distance=0.7cm]
\node[bdmnode] (u) {实验单元\\设备或班组};
\node[bdmnode,right=of u] (r) {随机分配\\按基线分层};
\node[bdmnode,right=of r] (o) {观测产品\\同组内可能相关};
\node[bdmnode,below=of r] (i) {按分配层级推断\\效应、区间、护栏};
\draw[bdmarrow] (u)--(r);
\draw[bdmarrow] (r)--(o);
\draw[bdmarrow] (o)|-(i);
\end{tikzpicture}
\caption{分配单元、观测单元与推断层级须同时声明。}
\end{figure}

共享设备或共享队列还可能造成跨组干扰：处理组提速改变了对照组等待。可考虑按产线、班次或时间块随机化。时间切换实验需要预设切换日历和清洗期\footnote{清洗期是切换处理后暂不计入效果比较的时段，用于减弱上一种处理的残留影响；其长度应依据业务机制和处理持续时间确定。}，并控制时段差异；存在不可忽略的干扰时，可定义产线或系统层面的策略效应。

\section{实验质量与持续监控}
\subsection{实验质量检查}

实验推断要求分配机制、观测资料和分析规则符合研究设计。分析前的质量检查至少包括以下内容。

\begin{table}[htbp]
\centering
\caption{实验质量检查}
\begin{tabular}{P{0.24\linewidth}P{0.34\linewidth}P{0.30\linewidth}}
\toprule
\rowcolor{BDMTableHead}
检查项 & 问题 & 风险\\
\midrule

随机分配 & 处理组和对照组是否按预设比例进入？ & 样本比例失衡可能意味着埋点或分流错误。\\
实验单元 & 分配单元和分析单元是否一致？ & 同一用户跨组会污染处理效果。\\
数据完整性 & 指标事件是否缺失或重复？ & 差异可能来自记录错误。\\
分析窗口 & 是否所有单元都有相同观察窗口？ & 后进入实验的单元可能被低估结果。\\
口径冻结 & 主指标和约束指标是否实验前定义？ & 事后挑选指标会放大偶然显著。\\
\bottomrule
\end{tabular}
\end{table}

样本比例失衡常被称为 SRM\footnote{SRM 是 sample ratio mismatch 的缩写，指实际入组比例与实验预设分配比例之间出现统计上难以由随机波动解释的偏离。}。若原计划 50/50 分流，但实际进入 B 组的样本显著偏少，随机分配、日志采集或资格筛选可能存在异常。此时需先核对随机分配、资格筛选和日志记录，再确定可用的分析样本。

\subsection{重复检验与多重比较}
\subjectindex{duochongbijiao}{多重比较}

A/B 测试在企业中常被频繁查看。反复检验并在首次显著时停止，会提高误报率，因为停止规则选择了产生显著结果的随机波动。反复查看结果时，需要采用序贯检验\footnote{序贯检验允许分阶段积累数据并判断是否停止，但检验边界须考虑多次查看结果带来的选择效应，才能维持规定的错误率。}或预先规定的中期分析规则；固定样本实验则应按事先设定的分析窗口和停止规则实施。

多重比较也会带来类似问题。若同时查看 20 个指标、10 个分群和多个时间窗口，即使没有真实效果，也很可能出现若干“显著”结果。实验报告应说明哪些指标是预先定义的主指标，哪些属于探索性分析。

\begin{remark}
探索性分组用于提出新的效应异质性假设。若某一客户组呈现较大效应，可以事先定义该组及其检验目标，在后续独立样本中开展确认性分析。
\end{remark}

\subsection{随时有效置信区间}
\subjectindex{zhixinxulie}{置信序列}
为构造同时适用于所有观察时点的区间，考虑独立同分布 \(X_i\in[0,1]\)、均值为 \(\mu\)。Hoeffding 不等式给出
\[
\Prob(\abs{\bar X_t-\mu}\ge r)\le2e^{-2tr^2}.
\]
将总错误预算分配为 \(\alpha_t=6\alpha/(\pi^2t^2)\)，有 \(\sum_{t\ge1}\alpha_t=\alpha\)。令
\[
r_t=\sqrt{\frac{\log(2/\alpha_t)}{2t}}
=\sqrt{\frac{\log(\pi^2t^2/(3\alpha))}{2t}}.
\]
对所有时间使用并集界\footnote{并集界指任意一组事件中至少一个发生的概率，不超过各事件概率之和；该上界不要求事件相互独立。}，
\[
\Prob\!\left(\exists t\ge1:
\mu\notin[\bar X_t-r_t,\bar X_t+r_t]\right)\le\alpha.
\]
将区间与 \([0,1]\) 取交集不改变覆盖保证。这是一个保守但可直接证明的置信序列。任意基于已观察数据的停止时点仍受到同一个错误概率控制；前提是独立、有界与共同均值等条件成立。一般自适应实验需要另外证明所需条件。

Howard 等人的工作发展了更紧的时间一致方法\cite{cs06}。比较两组时，可以为每组分配 \(\alpha/2\) 的错误概率预算，再由各自的均值区间构造差异区间。单组公式不能直接用于任意非独立的差分序列。

\subsection{效应、收益与推广条件}
若上线覆盖 \(M\) 个对象，每单位结果改善价值为 \(v\)，每个对象的实施成本为 \(c\)，固定投入为 \(F\)，则净收益估计为
\[
\widehat B=M(v\hat\tau-c)-F.
\]
它把统计效应转成业务阈值 \(\tau>(c+F/M)/v\)。一个置信区间完全大于零的效应，仍可能低于这个阈值；一个区间跨零的结果，也不能被解释为已经证明无效。推广还要求风险指标、目标人群与执行能力符合预设条件。

\section{企业实验与人工智能系统评价}

\subsection{Microsoft 的实验组织与运行}
\subjectindex{Microsoft}{Microsoft}
Microsoft Research 对在线实验平台的论文及可信实验实践说明了随机实验在产品迭代中的应用，并讨论组织与工程条件\cite{ms06}。CUPED 则把已有实验前信息转成统计效率。制造实验可借鉴其随机分配、质量检查和结果报告方法，并根据设备相关性及生产周期确定实验单元。

在制造场景的实验设计中，质量部门先确定最终良率与客户缺陷护栏，生产部门确定随机化和污染边界，数据团队冻结取数与结果窗口，决策者事先规定推广规则。实验记录应保存失败实验和负面结果，并分别报告事后探索与预先规定的检验。

\subsection{自然语言查询系统的比较实验}
以一个完整查询问题为评价单位。同一批问题分别提交给系统 A 与 B，记录是否正确完成、是否违反口径、延迟和总成本。令 \(S_{iA},S_{iB}\in\{0,1\}\)，成对差异为 \(D_i=S_{iB}-S_{iA}\)，估计改善量为 \(\bar D\)。相同任务的配对能控制题目难度，但任务之间若高度相似，应按任务族划分和分析。

运行评价分别测量 SQL 执行成功率、答案正确率和任务总用时。在线实验计入人工复核与返工时间，并记录模型、提示、语义库和工具的版本。数值结果可与确定性计算比对，开放性答案可由独立人员抽查，以减少生成与评价共享错误的影响。

\subsection{智能体参与运营决策的试点评价}
如果实验处理从“人工查询”改为“可调用查询与工单工具的智能体”，实验单位首先
应由业务流程确定。工厂同一工作中心的订单共享设备和人员，订单级随机化可能使
处理组加急挤占对照组产能。可以按相互独立的工作中心或时间段分配策略，并在
设计阶段考察订单跨组流动。任务层测试检验工具是否按规定工作；经营层随机实验
估计部署对交付结果的影响，两种证据应分别报告。

对组 \(g\) 的结果，可定义在固定观察窗口内的准时交付率 \(Y_g\)、总复核工时
\(H_g\) 和实际加急成本 \(C_g\)。主要结果可以是按到期订单数加权的交付率，
同时把错误修改交期、越权提交、客户投诉和人工返工作为预先规定的约束指标。
如果自动建议很快，却增加了主管复核时间，\(H_g\) 会呈现这项代价。
模型、数据快照、提示版本、工具权限和审批规则应在比较期间固定或按分配组
记录变化。对工作中心随机分配时，方差估计须以随机化单位为基础，不能把同一
工作中心的所有订单当成相互独立的实验单位。

在同一组封存订单上比较两种系统的任务完成率，可以检验给定输入下的正确性；
工作中心的试点则须规定处理分配与经营结果的观察窗口。若试点仅含两个工作中心且规模不足，
应着重分析设计的可行性与估计不确定性，不据此作显著性结论。

\section{实验结果与推广评价}

试点结果的推广涉及向目标总体外推。试点可能选在订单结构稳定、主管配合、
设备较新的产线，而全面推广会进入完全不同的环境。随机化消除了试点内部的系统性
分配差异，不自动消除试点与目标总体之间的差异。因而，推广报告应分别说明试点
内部比较是否可信，以及目标场景为何具有可比的执行条件。

若试点按用户随机，而推广后共享服务器负载显著上升，试点效果也可能不再保持。
制造中的共享设备、人员协作和物料竞争具有类似问题。此时要评价的是系统级策略
改变后的结果，而不只是固定环境下单个对象的处理效应。扩大试点范围、按适当集群
随机或显式建模相互作用，都是重新匹配估计目标与干预范围的办法。

实验报告应按研究设计、计算结果和推广评价组织，基本内容如下。
\begin{enumerate}
  \item 业务问题：本实验要支持什么管理决策；
  \item 实验设计：实验单元、处理、对照、随机化、纳入规则和分析窗口；
  \item 指标体系：主指标、约束指标、MDE、显著性水平和功效；
  \item 数据质量：样本量、分流比例、缺失、重复、异常和口径版本；
  \item 统计结果：效应估计、标准误、置信区间、p 值和样本量解释；
  \item 管理解释：收益、成本、风险、适用范围和不确定性；
  \item 行动建议：推广、继续实验、回滚、改版或分层复验；
  \item 复盘入口：实验代码、数据批次、报告版本和最终决策。
\end{enumerate}

实验报告保存设计、数据、分析结果和决策依据，便于复查当时的信息条件并规划后续研究。

\section{计算实例与结果解释}

本章样例数据位于：
\begin{center}
\path{data/sample/ab_test_checkout.csv}
\end{center}
该数据模拟一个结账页 A/B 测试。实验单元为访客，A 组为旧版结账页，B 组为新版结账页。主指标是是否转化，辅助指标是每访客收入，约束指标包括页面加载时间和客服工单率。

配套代码位于：
\begin{center}
\path{src/bdm_decision/cases/week06_ab_testing.py}
\end{center}

核心结果如下：
\begin{itemize}
  \item A 组样本量为 20，转化率为 0.3500，每访客收入为 16.95；
  \item B 组样本量为 20，转化率为 0.5500，每访客收入为 31.65；
  \item 转化率差异为 0.2000，p 值为 0.2059，95\% 置信区间跨越 0；
  \item 每访客收入差异为 14.70，95\% 置信区间跨越 0；
  \item 页面加载时间平均增加 51.65 毫秒，约束指标显示处理组可能恶化体验；
  \item 若基准转化率为 0.35，希望检测 10 个百分点的 MDE，在 80\% 功效下每组约需 376 个样本。
\end{itemize}

样例的转化率点估计较高，但区间较宽，同时加载时间增加。按要求正向效果明确且体验风险可接受的推广规则，当前结果尚未达到条件。后续可研究加载时间增加的原因，并根据预期效应、功效和观察期限设计新一轮实验。

\section*{本章小结}
\addcontentsline{toc}{section}{本章小结}

指标差异的总体解释需要统计推断。随机实验通过设计上的随机化减少选择偏差，差异估计量提供处理效果的自然估计，标准误和置信区间刻画估计不确定性，p 值提供原假设下的极端性摘要，功效和 MDE 限定实验所需的样本规模。推广评价综合考虑主指标、约束指标、样本代表性、业务成本及实验实施记录。

CUPED 通过实验前协变量降低估计方差，整群设计按照共同接受处理的单位安排比较，置信序列在所有观察时点共同控制覆盖错误。人工智能参与生产决策的试点应按实际执行单位安排随机化，同时记录交付结果、人员复核和运行成本。将效应估计用于推广时，还需考虑实施规模、成本、风险和目标人群。预测任务则研究在给定信息下如何估计未知结果，见第7章。

\chapterreading{随机化推断的经典来源包括 Fisher 与 Neyman\cite{fisher1935,neyman1923}，
实验安排可参考 Cox 以及 Box 等\cite{cox1958,box2005}。现代在线实验和组织治理
见 Kohavi 等\cite{kohavi2020}，现场实验见 Gerber 与 Green\cite{gerber2012}。
对 p 值的解释应结合 ASA 声明阅读\cite{wasserstein2016}。CUPED 和置信序列分别
提供降低方差与时间一致推断的方法\cite{cuped06,cs06}。应用这些方法仍需明确
实验单元、目标效应和错误控制。}

\section*{思考与习题}
\addcontentsline{toc}{section}{思考与习题}

\begin{exercise}[方差与设计]
从调整结果的方差展开式推导 CUPED 系数。若误用实验后设备温度作为协变量，为什么可能改变估计目标？在 10 台设备、每台 100 件的实验中，说明 1000 件为什么不等于 1000 个独立分配单元。
\end{exercise}
\begin{exercise}[序贯与业务阈值]
取 \(\alpha=0.05,t=1000\)，计算本章置信序列半径，并与普通固定样本区间比较。若 \(M=10000,v=20,c=0.4,F=2000\)，计算值得推广的最小效应。区间高于零是否已经足够？
\end{exercise}

\begin{exercise}
围绕 \path{data/sample/ab_test_checkout.csv} 完成下列任务：
\begin{enumerate}
  \item 分别计算 A 组和 B 组的样本量、转化率、每访客收入、平均订单价值、页面加载时间和客服工单率；
  \item 使用差异估计量计算 B 组相对 A 组的转化率提升，并给出 95\% 置信区间；
  \item 说明为什么“B 组转化率高 20 个百分点”不足以直接支持全量推广；
  \item 若管理层认为至少提升 5 个百分点才值得推广，请重新估算所需每组样本量；
  \item 检查该实验是否存在样本比例失衡、护栏恶化或指标口径风险；
  \item 为该实验写一段不超过 300 字的管理建议，必须同时提及主指标、约束指标和下一步行动。
\end{enumerate}
\end{exercise}
